\subsection{Interpretation and comparison with existing research}

This discussion answers the research questions in light of the six-dataset, four-operator benchmark, the branch-count ablation, the tuned-candidate checks, the mechanism evidence, and the synthetic multi-class scaling benchmark.

\subsubsection{RQ1: Does residual topology affect performance under a shared graph-classification protocol?}

The main benchmark argues against a universal residual topology. Across 24 dataset--operator combinations, MatrixRes is the most frequent winner and has the best mean rank, but all five model families win at least two combinations. Under GCNConv alone, Plain is strongest on PROTEINS, HorizontalRes is strongest on DD, MatrixResGated is strongest on ENZYMES, and MatrixRes is strongest on MUTAG, AIDS, and Mutagenicity. These outcomes indicate that the useful residual neighborhood depends on dataset-specific tolerance for branch specialization, residual traffic, and optimization complexity.

This conclusion is consistent with the broader GNN literature: changing the propagation operator, normalization, or pooling design often shifts the preferred architecture \cite{xu2019gin,hamilton2017inductive,velivckovic2017graph,ying2018hierarchical}. The present study adds that residual topology itself is another source of this shift. Matrix-style reuse can be read as a local residual stencil over the branch-by-layer grid, analogous in spirit to structured neighborhood mixing \cite{abu2019mixhop}, but applied to hidden-state reuse rather than to graph adjacency. Its advantage is therefore architectural organization, not a new message-passing primitive.

\textbf{Answer:} MatrixRes is the most frequent winner in the current full benchmark, but residual topology is still dataset- and operator-dependent. A graph-classification benchmark should compare vertical, horizontal, matrix-style, and sparse/gated reuse rather than assuming that one residual direction is universally best.

\subsubsection{RQ2: How do branch count, residual traffic, and mechanisms explain performance changes?}

The branch-count ablation shows that \(B\) is not only a capacity parameter. Increasing \(B\) adds branches, residual paths, and optimization interactions. In the GCNConv ablation, PROTEINS has a moderate useful region: HorizontalRes peaks at \(B=4\), MatrixRes at \(B=6\), and MatrixResGated at \(B=3\). DD prefers smaller budgets, with peaks at \(B=1\) or \(B=2\). The mechanism summaries are consistent with this split. On PROTEINS, MatrixRes with \(B=6\) produces the largest branch diversity, whereas HorizontalRes keeps lower branch cosine and a smaller residual footprint. On DD, all best branch-count settings have nearly saturated branch cosine and CKA, so extra residual traffic is less likely to create new useful functions. Branch diversity may increase with \(B\), but if branch similarity remains high or gradients weaken, the extra branches do not translate into better classification.

This rise-then-fall behavior is also a practical warning for residual GNN design. Residual connections are often introduced to make deeper or wider models easier to optimize \cite{he2016deep,li2019deepgcns,li2021training}, but a multi-branch residual model changes both capacity and routing density. A larger \(B\) increases the number of paths through which hidden states can be reinjected, so the useful regime is governed by the balance between functional diversity and residual interference. The branch-count results suggest that this balance differs sharply between PROTEINS and DD, even though both are commonly treated as protein-structure graph benchmarks.

\textbf{Answer:} In the PROTEINS/DD GCNConv ablation, branch count helps only within a limited operating region. Larger branch budgets are useful when they create functional diversity, and less useful when residual complexity grows faster than branch separation or optimization quality.

\subsubsection{RQ3: When is matrix-style residual reuse more effective on real data and controlled synthetic multi-class structures?}

MatrixResGated tests whether matrix residual traffic should be filtered before injection. The main benchmark shows that sparse/gated reuse is not uniformly better than dense matrix reuse: MatrixRes wins more dataset--operator combinations overall. However, MatrixResGated is still useful because it wins several combinations and provides a controlled way to reduce weak residual coefficients. The ENZYMES GCNConv win is reported but should be interpreted cautiously because all compared models have low absolute accuracy on that dataset. The mechanism table supports the residual-control interpretation: the gated variant generally keeps fewer active residual pathways than dense MatrixRes, while the learned gate remains in a moderate range rather than collapsing to zero. This distinction is important. The value of MatrixResGated is not that sparsity always improves accuracy, but that sparse residual control can be beneficial when matrix-style reuse would otherwise inject too much low-value residual traffic.

This places sparse/gated residual reuse between two familiar design ideas. Like residual gated graph networks, it learns how strongly historical information should be injected \cite{bresson2017residual}; like sparsification and regularization methods, it reduces the influence of weak or noisy pathways \cite{rong2019dropedge,cangea2018towards}. The present implementation is deliberately simple: it uses a scalar gate and soft-shrink filtering rather than a large routing network. This keeps the ablation interpretable and prevents the gated variant from becoming a separate high-capacity model.

MatrixResGated defines a tunable family rather than a plug-and-play winner. The fold-0 sensitivity scan suggests candidate operating regions: PROTEINS has a mild-sparsity candidate, while DD has conservative candidates based on smaller hidden dimension and lower initial gate. The five-fold candidate reruns refine this interpretation. On DD, the smaller hidden dimension remains beneficial and parameter-efficient. On PROTEINS, the mild-sparsity candidate does not preserve its fold-0 advantage across all folds.

\textbf{Real-data evidence.} Sparse/gated matrix reuse is a control mechanism for residual traffic. It is complementary to dense MatrixRes rather than a strict replacement, and candidate settings must be validated across folds before they are treated as performance claims. The current evidence supports a smaller controlled matrix model on DD and leaves PROTEINS sparsification as an unstable candidate rather than a confirmed gain.

The synthetic multi-class benchmark separates class granularity from real-dataset idiosyncrasies. In the two- and three-class settings, simpler residual choices remain competitive: Plain is best at \(C=2\), and HorizontalRes is best at \(C=3\). This result is important because it prevents an overbroad claim that matrix-style reuse is always needed. When the class count increases, the pattern changes. In the \(C=4,\ldots,8\) regime, MatrixResGated has the best average accuracy and normalized accuracy, while MatrixRes has the best macro-F1. MatrixRes also has the smallest accuracy drop from \(C=2\) to \(C=8\), suggesting that dense local matrix reuse is relatively stable as the label space becomes finer.

The distribution-level synthetic results refine this conclusion. MatrixRes is strongest in the higher-class BA, ER, and regular-degree control settings, where class labels are tied to degree growth, edge probability, or degree-uniform structural changes. These families share a common pattern: the label space is formed by related structural regimes rather than by unrelated categories, so local matrix residual reuse can combine depth-wise refinement with neighboring-branch exchange. SBM is community-driven and less favorable to matrix-style reuse, but MatrixRes remains close to the best result. Therefore, the synthetic benchmark supports a conditional claim: matrix-style residual reuse is most useful when the task requires separating several related structural regimes, while its advantage remains distribution-dependent.

\textbf{Answer:} Taken together, the real-data and synthetic scaling experiments support MatrixRes and MatrixResGated as conditionally effective rather than universal residual families. Sparse/gated reuse is useful for controlling excessive residual traffic, while matrix-style reuse becomes more favorable under finer structural discrimination, especially from four to eight classes. The effect is strongest on several controlled degree- and density-driven graph families, but it remains distribution-dependent.

\subsection{Limitations}

Several limitations bound the generality of these findings. First, although the main benchmark includes four message-passing operators, it still uses a fixed graph-level pooling and classification head; alternative readout designs may interact with residual topology. Second, the compared operators are canonical message-passing backbones rather than recent graph transformer or local-global hybrid architectures \cite{dwivedi2020generalization}. This is a deliberate control choice, but it means the paper should not be read as a state-of-the-art model comparison. Third, the sensitivity scan is fold-0 only; although selected candidates were rerun across five folds, that follow-up is limited to MatrixResGated on PROTEINS and DD. Fourth, ENZYMES has low absolute accuracy and high variability, so its result should be interpreted cautiously. Finally, the mechanism analysis is correlational. Branch diversity, cosine similarity, CKA, and gradient norms support a coherent explanation, but they do not prove causality.

Another limitation is that the benchmark intentionally avoids task-specific chemical or biological feature engineering. This makes the residual-topology comparison cleaner, but it also means that the absolute accuracies should not be interpreted as optimized state-of-the-art molecular prediction results. The synthetic benchmark has a related limitation: its labels are generated from controlled structural parameters, so it is useful for stress testing residual topology but cannot substitute for broader real multi-class graph benchmarks. The paper's contribution is methodological: it identifies when different residual pathways are helpful under matched conditions.

\subsection{Future directions}

Future work should extend the comparison to larger graph-classification datasets, richer readout mechanisms, recent local-global backbones, and adaptive residual-neighborhood selection. A natural next step is to replace manually chosen vertical, horizontal, or matrix stencils with a learned local residual policy. The tuned-candidate validation should also be broadened beyond the current MatrixResGated checks to determine whether dataset-specific residual control can be selected reliably.
